\chapter{Introduction}
\label{ch:intro}

Protein--protein binding affinity prediction asks how strongly two proteins bind. In benchmark datasets for this task, the target is a thermodynamic quantity such as dissociation constant $K_D$ or binding free energy $\Delta G$, while the input is usually a single three-dimensional complex structure \citep{Liu2024PPBAffinity}. Accurate affinity prediction is important for protein design, therapeutic targeting, and drug discovery \citep{Guo2022PPBAReview,Modell2016PPIs}. Experimental measurements provide the ground truth, but they are too slow and expensive for large exploratory searches. Physics-based free-energy calculations aim more directly at binding free energy, but the more faithful versions are also expensive \citep{Cournia2017RBFE,Siebenmorgen2020PPIReview}. Many practical machine-learning predictors therefore estimate affinity quickly from structural input alone.

One difficulty is that binding affinity is an ensemble property: it depends on the range of shapes and relative arrangements the proteins sample across an ensemble, beyond any one frozen structure. Proteins can bend, shift, and sometimes slide against each other as they bind \citep{HenzlerWildman2007Dynamic}, so a single bound complex can miss relevant information. Fast structure-based predictors then face two linked design choices. Some recent protein--protein models focus tightly on the interface, while others use whole-complex or hybrid representations \citep{Liu2024PPBAffinity,Moal2011AffinityPrediction,RomeroMolina2022PPIAffinity,Xie2025PPIGraphomer}. Another design choice, explored more extensively so far in protein--ligand affinity prediction, is whether MD can add useful information beyond one bound structure \citep{Min2024Dynaformer,Libouban2025Spatiotemporal}. In this setting, proteins are represented as residue graphs, with amino-acid residues as nodes and residue--residue proximity or contact as edges, and the comparison tests both choices in one controlled protein--protein setting.

The analysis uses a curated wild-type, that is, non-mutant, subset of PPB-Affinity \citep{Liu2024PPBAffinity}. The final comparison uses a matched subset of 1437 complexes from a curated table of 2818. Each retained complex went through the same preparation and short-MD procedure before conversion into a post-preparation residue graph with DeepRank2 and enrichment with trajectory-derived summaries. Affinity was then predicted in two ways: a masked GraphVAE supplied latent means for ridge regression, and a direct supervised graph regressor predicted affinity directly from the graph. The shared setup tests whether short MD summaries improve on a static baseline.

The comparison asks four questions. Is a full residue graph better than an interface-restricted representation? Do MD-derived summaries improve prediction over static graph features alone? Does a masked GraphVAE produce a better predictive representation than a direct supervised graph regressor with the same encoder architecture? How stable are those answers under repeated held-out evaluation?

Chapter~\ref{ch:litreview} places the work in the literature on protein--protein affinity prediction, molecular simulation, and learned structural representations. Chapter~\ref{chap:methods} describes the methods. Chapter~\ref{ch:results} reports the 52-model comparison. Chapter~\ref{ch:discussion} gives the main conclusions. The appendices cover archived approaches, the complete model table, and the repository layout.

The implementation is organised in \url{https://github.com/RohanDevereux/Bridging}. Appendix~\ref{ch:appendix_archive} records the alternative approaches explored before the final study design was settled.
